\documentclass[11pt]{article}
\title{The Maximum Number of Plane Regions Cut by Straight Lines: A Sharp Incremental Bound}
\date{}
\author{SIGMA AI-assisted research synthesis}
\usepackage[T1]{fontenc}
\usepackage[utf8]{inputenc}
\usepackage[margin=1in]{geometry}
\usepackage{amsmath,amssymb,booktabs,longtable}
\usepackage{microtype}
\usepackage{titlesec}
\emergencystretch=2em
\raggedbottom
\linespread{1.07}
\usepackage{xcolor}
\definecolor{blue}{HTML}{24485E}
\titleformat{\section}{\Large\bfseries\sffamily\color{blue}}{\thesection}{0.8em}{}
\titleformat{\subsection}{\large\bfseries\sffamily}{\thesubsection}{0.7em}{}
\setlength{\parskip}{.6em}
\setlength{\parindent}{0pt}
\begin{document}
\maketitle
\begin{center}\small\textit{AI-generated research synthesis \textperiodcentered{} SIGMA}\end{center}

\begin{abstract}
How many connected regions can a finite collection of full straight lines create in the Euclidean plane? For every nonnegative integer $n$, the maximum is $M(n)=1+n(n+1)/2$. The argument rests on an exact increment rule: a new line meeting the existing arrangement at $m$ distinct points creates precisely $m+1$ additional regions, unless it coincides with an existing line. Convexity of the existing regions makes this local count valid. The rule gives a universal upper bound, including arrangements with parallel, concurrent, or repeated lines. The explicit family $y=ix+i^2$ attains every incremental bound and therefore realizes the maximum for every $n$. The empty arrangement and small cases complete the boundary checks.

This provides a self-contained deductive argument, without a claim of independent formal verification.
\end{abstract}
\section*{About this paper}
This paper was generated by SIGMA, an AI-assisted research system, from the project's recorded investigations, findings, and supporting evidence.
The analysis and conclusions reflect an automated synthesis of that research. Although SIGMA incorporates evidence tracking, critical review, and consistency checks, the resulting paper has not necessarily undergone independent verification or peer review.
Readers should evaluate its arguments, supporting evidence, and stated limitations before relying on its conclusions.
Research style: Distinguish proved statements from conjectures, failed arguments, and open proof obligations. State assumptions and explain bridges between logical steps.
\textit{Claims discussed: }C2, C3, C4, C5, C6, C7, C9
\section{Introduction}
Let $n$ be a nonnegative integer, and consider $n$ full straight lines in the Euclidean plane. A region is a connected component of the complement of their union. Write $M(n)$ for the largest possible number of regions; the task is to bound all arrangements and exhibit an arrangement attaining the bound for each $n$.

The upper-bound argument permits parallel lines, several lines through one point, and repeated lines. These possibilities must not be excluded by an unstated position assumption. A sharp result therefore requires two separate arguments: a universal restriction on region creation and an admissible construction that reaches that restriction.

\textit{Claims discussed: }C1, C8
\section{Evidence and scope}
The analysis counts regions by inserting lines successively. It first identifies regions as nonempty intersections of open half-planes, then uses their convexity to establish an exact rule for the increase caused by one new line. Summing the resulting incremental bounds gives the universal upper bound.

Attainment is analyzed separately through the explicit family $L_i:y=ix+i^2$. Solving the pairwise intersection equations shows that each new line meets all earlier lines at distinct points, which is exactly the condition needed to reach its incremental bound. Direct substitution and the corresponding geometric increments check the cases $n=0,1,2,3,4$; these checks do not replace the general argument.

\textit{Claims discussed: }C2, C3, C4, C5, C7, C8
\section{Lemma 1. Regions are open and convex}
Choose an affine equation for each line. On any connected component of the complement, each equation has a constant strict sign: a continuous function cannot take both positive and negative values on a connected set while avoiding zero. Thus each region lies in a set specified by choosing one open half-plane for every line.

Conversely, every nonempty set specified by such strict signs is an intersection of open half-planes. It is open and convex, hence connected, and it avoids every line. A component is contained in one such set, and that connected set is contained in the same component, so the two coincide.

\textit{Claims discussed: }C1, C2
\section{Lemma 2. The exact increase from one new line}
Let $L$ be a new line that does not coincide with an existing line, and let $m$ count its distinct intersection points with the existing arrangement. Removing those points leaves exactly $m+1$ connected pieces: intervals and rays, or the whole line when $m=0$. Each piece avoids all old lines and therefore lies in a single old region.

Different pieces lie in different old regions. Otherwise, points from two pieces would belong to the same convex region, so the segment joining them would also lie in that region. That segment contains an old-line intersection point, contradicting the definition of a region.

Every old region met by $L$ is split into exactly two regions. A meeting point has a small disk contained in the old region, so the region has points on both sides of $L$. Its intersections with the two open half-planes bounded by $L$ are nonempty and convex; they are connected and separated by the removed line.

Regions not met by $L$ remain unchanged. Consequently, the $m+1$ pieces each account for one additional region, giving an exact increase of $m+1$. If $L$ coincides with an old line, the union of lines does not change and the increase is zero.

\textit{Claims discussed: }C2, C3
\section{Proposition. A universal upper bound}
Insert an arbitrary arrangement in any order, and let $r_k$ denote its region count after the first $k$ lines. A noncoincident $k$th line meets each of the preceding $k-1$ lines at most once, so it has at most $k-1$ distinct intersection points. Lemma 2 therefore gives $r_k-r_{k-1}\leq k$; a repeated line satisfies the same inequality because it adds nothing.

With no lines, the whole plane is one region, so $r_0=1$. Summing the incremental inequalities yields

$r_n\leq 1+\sum_{k=1}^{n}k=1+\frac{n(n+1)}{2}.$

Parallel lines and several lines meeting at one point cannot invalidate this estimate: they reduce, rather than increase, the number of distinct intersections available on a new line.

\textit{Claims discussed: }C3, C4
\section{Construction. An arrangement attaining every bound}
For each positive integer $i$, define $L_i$ by $y=ix+i^2$. The slopes are distinct, so these lines are neither parallel nor coincident. For $i<k$, equating the two line equations gives

$ix+i^2=kx+k^2,\qquad x=-(i+k),\qquad y=-ik.$

For fixed $k$, the coordinates $-(i+k)$ are distinct as $i$ ranges from $1$ to $k-1$. Hence $L_k$ meets the preceding lines at exactly $k-1$ distinct points, and Lemma 2 shows that it adds exactly $k$ regions. This also covers $k=1$: the first line has no old intersection points and adds one region.

Taking the first $n$ lines therefore gives exactly

$1+\sum_{k=1}^{n}k=1+\frac{n(n+1)}{2}$

regions. The construction works for every positive integer $n$, so attainment does not rely on an unspecified generic placement.

\textit{Claims discussed: }C3, C5
\section{Theorem and boundary checks}
For every nonnegative integer $n$, the maximum number of regions formed by $n$ full straight lines in the Euclidean plane is

$M(n)=1+\frac{n(n+1)}{2}.$

The universal bound proves that no arrangement exceeds this number, and the explicit construction proves that the bound is attained for positive $n$. For $n=0$, the empty arrangement leaves one region, agreeing with the formula and the empty sum.

For $n=0,1,2,3,4$, the formula gives $1,2,4,7,11$. Geometrically, these correspond to the empty plane, one dividing line, two intersecting lines, three pairwise-intersecting nonconcurrent lines, and a fourth line crossing the previous three at distinct points. The successive increases are $1,2,3,4$, as required by the construction.

\textit{Claims discussed: }C4, C5, C6, C7
\section{Verification status and limitations}
\begin{itemize}
\item The argument applies to full straight lines in the Euclidean plane and to connected components of their complement. No result for segments, curves, or a bounded observation window is claimed.
\item The proof is supported by explicit deductions, not by independent certification. The supplied record contains no external formal verification, and agreement within its detailed review is not a substitute for such verification (C9).
\end{itemize}
\section{Open obligations}
\begin{itemize}
\item Independent verification of the complete argument remains outstanding. No specific unresolved mathematical inference is identified in the supplied detailed review.
\end{itemize}
\section{Frozen obligation reconciliation}
Unlisted obligations are not presumed resolved. The archived research snapshot retains complete source references and the full reconciliation.
\begin{itemize}
\item Convexity and identification of arrangement regions: established by reasoning step 1 of Investigator step 5ae10b4f-29aa-4770-91a8-aa7d378719e3, cycle 1. \textit{ (Not listed by author; not presumed resolved; 2 saved source reference(s))}
\item Exact increment m+1 for a noncoincident line and zero for a coincident line: established by reasoning step 2 of Investigator step 5ae10b4f-29aa-4770-91a8-aa7d378719e3, cycle 1. \textit{ (Not listed by author; not presumed resolved; 2 saved source reference(s))}
\item Universal upper bound for arbitrary arrangements: established by reasoning step 3 of Investigator step 5ae10b4f-29aa-4770-91a8-aa7d378719e3, cycle 1. \textit{ (Not listed by author; not presumed resolved; 2 saved source reference(s))}
\item Existence of an attaining arrangement for every positive integer n: established by the intersection calculation and increment argument in reasoning steps 4–5 of Investigator step 5ae10b4f-29aa-4770-91a8-aa7d378719e3, cycle 1. \textit{ (Not listed by author; not presumed resolved; 2 saved source reference(s))}
\item Empty arrangement and small-value checks: established by reasoning steps 5–6 of Investigator step 5ae10b4f-29aa-4770-91a8-aa7d378719e3, cycle 1. \textit{ (Not listed by author; not presumed resolved; 2 saved source reference(s))}
\item Independent verification of the complete argument remains outstanding. No specific unresolved mathematical inference is identified in the supplied detailed review. \textit{ (Author-listed; no exact frozen-record match; 0 saved source reference(s))}
\end{itemize}
\section{Conclusion}
The sharp maximum is $M(n)=1+n(n+1)/2$ for every nonnegative integer $n$. Convexity justifies the exact local region count, the intersection bound limits each insertion, and the family $y=ix+i^2$ realizes every allowed increase. The empty arrangement and the small-value checks agree with the general result.

\textit{Claims discussed: }C2, C3, C4, C5, C6, C7
\section*{Claim status and evidence}
\subsection{C1 -- Established}
The problem concerns finitely many full straight lines in the Euclidean plane, with regions defined as connected components of their complement; the universal bound allows parallel, concurrent, and repeated lines.
These are the stated domain and definitions of the saved argument, not additional geometric hypotheses.
\textit{Evidence: }project:01af8ffa-60c3-4eea-bc15-36dca347ddfa, artifact:f60b98af-2fd0-493d-8eb9-1b7b204be8fa:4b50ca5d-ad1f-4f6a-8042-81975427938f
\subsection{C2 -- Established}
Every region of a finite line arrangement is an open convex set, namely a nonempty intersection of open half-planes determined by strict signs of affine line equations.
Each affine equation has constant strict sign on a connected component avoiding its zero line. Conversely, every nonempty set with a prescribed collection of strict signs is convex and connected, so it is precisely a component. This is the first saved deductive step.
\textit{Evidence: }artifact:f60b98af-2fd0-493d-8eb9-1b7b204be8fa:4b50ca5d-ad1f-4f6a-8042-81975427938f
\subsection{C3 -- Established}
A new line not coinciding with an existing line adds exactly $m+1$ regions when it meets the existing lines at $m$ distinct points; a coincident line adds none.
The intersection points leave $m+1$ connected pieces on the new line. Convexity places these pieces in different old regions. Openness and convexity ensure that each region met by the line becomes exactly two regions, while unmet regions remain unchanged.
\textit{Evidence: }artifact:f60b98af-2fd0-493d-8eb9-1b7b204be8fa:4b50ca5d-ad1f-4f6a-8042-81975427938f
\subsection{C4 -- Established}
Every arrangement of $n$ lines has at most $1+n(n+1)/2$ regions.
The $k$th noncoincident line has at most $k-1$ distinct intersection points with earlier lines and therefore adds at most $k$ regions. A repeated line adds zero. Summing from the initial single region gives the bound.
\textit{Evidence: }artifact:f60b98af-2fd0-493d-8eb9-1b7b204be8fa:4b50ca5d-ad1f-4f6a-8042-81975427938f
\subsection{C5 -- Established}
For every positive integer $n$, the lines $L_i:y=ix+i^2$, with $1\leq i\leq n$, attain every incremental upper bound.
Distinct slopes prevent parallel or coincident pairs. For $i<k$, the intersection is $(-(i+k),-ik)$. With $k$ fixed, these points have distinct horizontal coordinates, so the $k$th line has exactly $k-1$ distinct intersections and adds exactly $k$ regions.
\textit{Evidence: }artifact:f60b98af-2fd0-493d-8eb9-1b7b204be8fa:4b50ca5d-ad1f-4f6a-8042-81975427938f
\subsection{C6 -- Established}
For every nonnegative integer $n$, the maximum is $M(n)=1+n(n+1)/2$, attained by the explicit family for positive $n$ and by the empty arrangement for $n=0$.
The universal upper bound and the construction's exact count coincide. For $n=0$, the empty union leaves the entire plane connected and the empty sum contributes zero.
\textit{Evidence: }artifact:f60b98af-2fd0-493d-8eb9-1b7b204be8fa:4b50ca5d-ad1f-4f6a-8042-81975427938f
\subsection{C7 -- Established}
The maximum counts for $n=0,1,2,3,4$ are respectively $1,2,4,7,11$.
These values result from direct substitution and agree with the construction's successive increases of $1,2,3,4$ from the initial single region.
\textit{Evidence: }artifact:f60b98af-2fd0-493d-8eb9-1b7b204be8fa:4b50ca5d-ad1f-4f6a-8042-81975427938f
\subsection{C8 -- Established}
The saved analysis is deductive: it separates the universal incremental bound from the attaining construction and treats small values only as checks.
The saved method and its six reasoning steps carry out precisely this separation. No computational result is used as a premise.
\textit{Evidence: }strategy:ee7d3b31-9c8d-4b44-a1a8-b123c67e59df, artifact:f60b98af-2fd0-493d-8eb9-1b7b204be8fa:4b50ca5d-ad1f-4f6a-8042-81975427938f
\subsection{C9 -- Limitation}
The supplied material supports a complete deductive argument but does not establish independent or formal verification; its aggregate completion labels conflict with its detailed mathematical review.
The detailed review accepts all six deductions, reports the objective achieved, and lists no remaining mathematical obligations, while calling the proof convincing but unverified. The aggregate outcome instead labels the result partial and incomplete without identifying a missing inference. The established labels here refer to the displayed deductions, not to independent certification.
\textit{Evidence: }session:f60b98af-2fd0-493d-8eb9-1b7b204be8fa, artifact:f60b98af-2fd0-493d-8eb9-1b7b204be8fa:4b50ca5d-ad1f-4f6a-8042-81975427938f
\section*{Evidence and authorship note}
Snapshot SHA-256: 56e843297cd316019f170c929023fd090bf8ba9772274635c121e66a40f9c3a9
Generated by an AI author from the cited saved research snapshot. This document is not independent verification. No external bibliography was added.
\section*{Validation warnings}
\begin{itemize}
\item The evidence has inconsistent completion labels: the aggregate outcome calls the argument incomplete and claims no solution, whereas the detailed review accepts every deductive step, reports the objective achieved, and records no remaining obligations. No concrete missing proof step accompanies the aggregate label.
\item The aggregate metrics describe five items as open proof obligations, but their texts explicitly say they are established, and the detailed remaining-obligation lists are empty. They have not been treated here as five unresolved mathematical gaps.
\item Claims marked established are supported by the reproduced deductive argument, not by agent agreement or formal verification. Independent verification remains outstanding.
\item No unsupported mathematical extension, computational experiment, external citation, or omitted material proof step was identified in preparing this draft.
\item C1 is labeled established, but cited saved evidence includes unverified, incomplete, or conjectural material.
\item C1 is labeled established while the frozen evidence contains counterexample or contradiction language; inspect the cited sources before accepting it.
\item C2 is labeled established, but cited saved evidence includes unverified, incomplete, or conjectural material.
\item C2 is labeled established while the frozen evidence contains counterexample or contradiction language; inspect the cited sources before accepting it.
\item C3 is labeled established, but cited saved evidence includes unverified, incomplete, or conjectural material.
\item C3 is labeled established while the frozen evidence contains counterexample or contradiction language; inspect the cited sources before accepting it.
\item C4 is labeled established, but cited saved evidence includes unverified, incomplete, or conjectural material.
\item C4 is labeled established while the frozen evidence contains counterexample or contradiction language; inspect the cited sources before accepting it.
\item C5 is labeled established, but cited saved evidence includes unverified, incomplete, or conjectural material.
\item C5 is labeled established while the frozen evidence contains counterexample or contradiction language; inspect the cited sources before accepting it.
\item C6 is labeled established, but cited saved evidence includes unverified, incomplete, or conjectural material.
\item C6 is labeled established while the frozen evidence contains counterexample or contradiction language; inspect the cited sources before accepting it.
\item C7 is labeled established, but cited saved evidence includes unverified, incomplete, or conjectural material.
\item C7 is labeled established while the frozen evidence contains counterexample or contradiction language; inspect the cited sources before accepting it.
\item C8 is labeled established, but cited saved evidence includes unverified, incomplete, or conjectural material.
\item C8 is labeled established while the frozen evidence contains counterexample or contradiction language; inspect the cited sources before accepting it.
\end{itemize}
\end{document}
